\section{Lexical Analysis}\label{c:lexical}

\subsection{Scanner and lexer design}

The lexical layer is written from scratch --- no parser generator, no \texttt{lex} --- and
runs in two passes. \textbf{Scanning} splits the input into typed character runs (word,
number, punctuation, separator, space, newline) using the five regular expressions in
Table~\ref{tab:regex}; each run keeps its byte offsets so every later artefact can point
back at the exact source span (\req{FR-10}). \textbf{Lexing} consumes those runs and
resolves each word against the lexicon, producing tokens carrying a terminal class,
canonical form, donor-language label and gloss (\req{FR-11}, \req{FR-12}).

\begin{table}[htbp]
\centering
\caption{Regular expressions driving the scanner, exported from \texttt{core/text.py}.}
\label{tab:regex}
\begin{tabular}{@{}l p{0.40\linewidth} p{0.32\linewidth}@{}}
\toprule
\textbf{Span kind} & \textbf{Regular expression} & \textbf{Matches} \\
\midrule
\texttt{WORD} & \ttfamily\scriptsize (\allowbreak{}?\allowbreak{}:\allowbreak{}[\allowbreak{}\textasciicircum{}\allowbreak{}\textbackslash{}\allowbreak{}W\allowbreak{}\textbackslash{}\allowbreak{}d\allowbreak{}\_\allowbreak{}]\allowbreak{}|\allowbreak{}[\allowbreak{}\textbackslash{}\allowbreak{}u\allowbreak{}0\allowbreak{}3\allowbreak{}0\allowbreak{}0\allowbreak{}-\allowbreak{}\textbackslash{}\allowbreak{}u\allowbreak{}0\allowbreak{}3\allowbreak{}6\allowbreak{}f\allowbreak{}]\allowbreak{})\allowbreak{}+\allowbreak{}(\allowbreak{}?\allowbreak{}:\allowbreak{}[\allowbreak{}'\allowbreak{}\textbackslash{}\allowbreak{}u\allowbreak{}2\allowbreak{}0\allowbreak{}1\allowbreak{}9\allowbreak{}-\allowbreak{}]\allowbreak{}(\allowbreak{}?\allowbreak{}:\allowbreak{}[\allowbreak{}\textasciicircum{}\allowbreak{}\textbackslash{}\allowbreak{}W\allowbreak{}\textbackslash{}\allowbreak{}d\allowbreak{}\_\allowbreak{}]\allowbreak{}|\allowbreak{}[\allowbreak{}\textbackslash{}\allowbreak{}u\allowbreak{}0\allowbreak{}3\allowbreak{}0\allowbreak{}0\allowbreak{}-\allowbreak{}\textbackslash{}\allowbreak{}u\allowbreak{}0\allowbreak{}3\allowbreak{}6\allowbreak{}f\allowbreak{}]\allowbreak{})\allowbreak{}+\allowbreak{})\allowbreak{}* & Letters and combining accents, with internal apostrophes or hyphens \\
\addlinespace
\texttt{NUMBER} & \ttfamily\scriptsize \textbackslash{}\allowbreak{}d\allowbreak{}+ & A run of digits \\
\addlinespace
\texttt{SPACE} & \ttfamily\scriptsize [\allowbreak{}\textasciicircum{}\allowbreak{}\textbackslash{}\allowbreak{}S\allowbreak{}\textbackslash{}\allowbreak{}r\allowbreak{}\textbackslash{}\allowbreak{}n\allowbreak{}]\allowbreak{}+ & Horizontal whitespace only \\
\addlinespace
\texttt{NEWLINE} & \ttfamily\scriptsize [\allowbreak{}\textbackslash{}\allowbreak{}r\allowbreak{}\textbackslash{}\allowbreak{}n\allowbreak{}]\allowbreak{}+ & One or more line breaks \\
\addlinespace
\texttt{PUNCTUATION} & \ttfamily\scriptsize [\allowbreak{}.\allowbreak{},\allowbreak{}!\allowbreak{}?\allowbreak{};\allowbreak{}:\allowbreak{}\textbackslash{}\allowbreak{}"\allowbreak{}(\allowbreak{})\allowbreak{}«\allowbreak{}»\allowbreak{}\textbackslash{}\allowbreak{}u\allowbreak{}2\allowbreak{}0\allowbreak{}1\allowbreak{}8\allowbreak{}\textbackslash{}\allowbreak{}u\allowbreak{}2\allowbreak{}0\allowbreak{}1\allowbreak{}9\allowbreak{}\textbackslash{}\allowbreak{}u\allowbreak{}2\allowbreak{}0\allowbreak{}1\allowbreak{}c\allowbreak{}\textbackslash{}\allowbreak{}u\allowbreak{}2\allowbreak{}0\allowbreak{}1\allowbreak{}d\allowbreak{}…\allowbreak{}]\allowbreak{}+ & Sentence and clause marks, including guillemets \\
\addlinespace
\bottomrule
\end{tabular}
\end{table}


The word regex is the one doing real linguistic work:

\begin{lstlisting}[language=Python, numbers=none, basicstyle=\ttfamily\footnotesize]
_WORD_RE = (?:[^\W\d_]|[\u0300-\u036f])+
           (?:['\u2019-](?:[^\W\d_]|[\u0300-\u036f])+)*
\end{lstlisting}

\begin{itemize}[leftmargin=1.4em, itemsep=0.05em]
  \item \term{[\^{}\textbackslash W\textbackslash d\_]} matches any Unicode letter while
        excluding digits and underscore, so \fw{ça}, \fw{où} and \fw{tchop} scan as single
        words without an explicit accent list.
  \item \term{[\textbackslash u0300-\textbackslash u036f]} admits \emph{combining} marks,
        so a decomposed \fw{e + U+0301} scans identically to a precomposed \fw{é} --- the two
        look identical on screen but are different byte sequences.
  \item The trailing group keeps elisions and hyphenated compounds whole: \fw{j'ai},
        \fw{n'est}, \fw{est-ce}. Both the ASCII apostrophe and the typographic \fw{U+2019}
        are accepted, because phones produce the latter.
\end{itemize}

\subsection{Token categories}

\begin{center}\begin{tabular}{@{}lrrl@{}}
\toprule
\textbf{Terminal} & \textbf{Entries} & \textbf{Slang} & \textbf{Example} \\
\midrule
VERB & 290 & 97 & \texttt{go} \\
NOUN & 210 & 192 & \texttt{combi} \\
ADV & 35 & 13 & \texttt{même} \\
ADJ & 33 & 23 & \texttt{cher} \\
DET & 30 & 1 & \texttt{les} \\
PRON & 23 & 1 & \texttt{on} \\
PREP & 22 & 0 & \texttt{au} \\
INTJ & 19 & 19 & \texttt{ashia} \\
CONJ & 14 & 0 & \texttt{qu'il} \\
FORMULA & 7 & 7 & \texttt{on dit quoi} \\
NEG & 5 & 0 & \texttt{ne} \\
\midrule
\textbf{Total} & \textbf{688} & \textbf{353} & \\
\bottomrule
\end{tabular}
\end{center}

The \TerminalCount\ terminals cover the categories the assignment asks for --- nouns, verbs,
slang and code-mixed expressions --- plus three the data forced on us. \term{NEG} is split
from \term{ADV} because Francanglais negation is positionally constrained in a way ordinary
adverbs are not, and folding the two together introduced grammar conflicts. \term{INTJ}
covers the discourse particles (\fw{ehh}, \fw{ashia}, \fw{wah}) that open and close a large
share of real utterances. \term{FORMULA} covers lexically frozen multiword expressions.
Additionally, \term{SEP} is emitted for \term{,} \term{;} \term{:}, which are genuine clause
boundaries in speech; sentence-final \term{.\;!\;?} remain trivia.

\subsection{Slang and code-mixed tokens}

Every entry records its donor language, which is what lets the analyzer report code-mixing
quantitatively rather than impressionistically. A statement is reported as code-mixed when
its tokens resolve to more than one donor label, and the evidence reported is the list of
tokens and labels themselves.

\begin{center}\begin{tabular}{@{}lr@{}}
\toprule
\textbf{Donor language} & \textbf{Token entries} \\
\midrule
french & 425 \\
uncertain & 132 \\
english & 50 \\
pidgin & 23 \\
mixed & 14 \\
duala & 14 \\
beti & 11 \\
bassa & 6 \\
bamileke & 5 \\
fulfulde & 3 \\
hausa & 2 \\
bulu & 1 \\
german & 1 \\
lingala & 1 \\
spanish & 1 \\
\bottomrule
\end{tabular}
\end{center}

Slang is not a separate terminal --- a slang noun is still a \term{NOUN} --- because slang is a
fact about \emph{provenance}, not about syntax. Treating it otherwise would have duplicated
the whole grammar.

\subsection{Accent folding and homograph protection}

Two normalisations are applied, and the distinction matters.
\texttt{canonicalize\_word} (NFC normalise, then casefold) is the identity used for storage
and display, so \fw{Taxi}, \fw{TAXI} and \fw{taxi} are one entry.
\texttt{fold\_word} (NFD decompose, strip combining marks, recompose, casefold) is the
\emph{fallback} identity, so \fw{bénéf} and \fw{benef} reach the same entry.

Folding is deliberately a fallback rather than the primary key, because in French it
destroys real contrasts. Lookup proceeds in strict priority order: exact multiword, folded
multiword, exact single-word, folded single-word, otherwise \term{UNKNOWN}. Because exact
matches always win, the four attested homograph pairs survive intact (\req{FR-15}):

\begin{center}
\small
\begin{tabular}{@{}llll@{}}
\toprule
\textbf{Form A} & \textbf{Class} & \textbf{Form B} & \textbf{Class} \\
\midrule
\fw{a}    & \term{VERB} (\textit{has})  & \fw{à}    & \term{PREP} (\textit{to, at}) \\
\fw{la}   & \term{DET} (\textit{the})   & \fw{là}   & \term{ADV} (\textit{there}) \\
\fw{ou}   & \term{CONJ} (\textit{or})   & \fw{où}   & \term{ADV} (\textit{where}) \\
\fw{mais} & \term{CONJ} (\textit{but})  & \fw{maïs} & \term{NOUN} (\textit{maize}) \\
\bottomrule
\end{tabular}
\end{center}

When a folded lookup resolves to a form that has a homograph, the analyzer emits a hint
naming both readings instead of choosing one silently.

\subsection{Multiword expressions and unknown words}

\LexiconMultiword\ entries span more than one word (\fw{est-ce que}, \fw{on va faire
comment}). The lexer indexes them by first word, attempts the longest match first, and
refuses to match across a punctuation or separator boundary, so a comma between two words
prevents them from being glued into one token (\req{FR-14}).

An unrecognised word is not an error and is not discarded (\req{FR-16}). It is emitted as an
\term{UNKNOWN} token with its source span intact, which produces three useful effects: the
parser fails at a precise location rather than vaguely; the suggester can offer near matches
computed over \emph{folded} forms, so an accent slip still finds its target; and the reviewer
gets a concrete list of vocabulary the lexicon is missing.
